% ============================================================
% Capítulo 5 — Metodología
% ============================================================

\chapter{Metodología}
\label{cap:metodologia}

Este capítulo describe el protocolo común a todos los experimentos: cómo
se parten los datos, cómo se entrena y se elige cada modelo, cómo se
obtienen las instancias y cómo se evalúan. Todas las ejecuciones siguen
las mismas reglas, fijadas antes de empezar a comparar. Así, cualquier
diferencia entre dos resultados se debe a la pieza que cambia (la
pérdida, el post-proceso o los hiperparámetros) y no al protocolo. El
código está disponible en \url{https://github.com/AgustinG-git/tfm-filamentos}.

% ------------------------------------------------------------
\section{Datos y partición}
\label{sec:particion}
% ------------------------------------------------------------

Se usa el conjunto de entrenamiento de la competición, con las
anotaciones de MAGFiLO~\cite{ahmadzadeh2024,kaggle2026}. El conjunto de
test no tiene anotaciones públicas y solo se usa para enviar
predicciones al \emph{leaderboard}.

\paragraph{Partición por días.} Las imágenes de un mismo día son casi
idénticas: el Sol apenas cambia en unas horas. Si una imagen de un día
estuviera en entrenamiento y otra del mismo día en validación, el modelo
se evaluaría sobre algo que ya ha visto. Por eso la partición se hace por
\emph{fecha}: ningún día aparece en dos subconjuntos. Se reserva el
20\,\% de los días para validación, y estos se dividen a su vez en dos
mitades, $A$ y $B$ (Tabla~\ref{tab:split}). La partición se genera una
sola vez, con semilla fija, y se guarda junto al código.

\begin{table}[ht]
\centering
\begin{tabular}{@{}lrrr@{}}
\toprule
Subconjunto & Días & Imágenes & Pares anotador--imagen\\
\midrule
Entrenamiento & 524 & 561 & 920\\
Validación $A$ & 66 & 72 & 113\\
Validación $B$ & 66 & 74 & 121\\
\bottomrule
\end{tabular}
\caption{Partición de los datos. La unidad de evaluación es el par
anotador--imagen (Sección~\ref{sec:competicion}).}
\label{tab:split}
\end{table}

\paragraph{Ajuste en $A$, reporte en $B$.} Toda decisión se toma en
$A$: la época del modelo, los parámetros del post-proceso y los
hiperparámetros. La cifra que se reporta es la de $B$, que no interviene
en ninguna decisión. Si se eligiera en el mismo conjunto en el que se
mide, la cifra estaría inflada: entre muchas variantes, alguna sale
mejor por azar.

\paragraph{Preprocesado.} Las imágenes H-alfa son de un canal y
$2048\times2048$ píxeles. Para entrenar se reducen a $512\times512$
(interpolación por área), se replica el canal tres veces y se
normalizan con la media y la desviación de ImageNet~\cite{deng2009},
porque el codificador viene preentrenado en ese conjunto. No se usan
aumentos de datos.

\paragraph{Etiquetas.} Cada anotación es un ejemplo de entrenamiento
independiente: una imagen con tres anotadores aporta tres pares
(imagen, máscara). La máscara de un par es la unión de los filamentos de
ese anotador, reducida a $512\times512$. La referencia de evaluación,
en cambio, son siempre las anotaciones originales, una por una y a
resolución completa.

% ------------------------------------------------------------
\section{Modelo y entrenamiento}
\label{sec:entrenamiento}
% ------------------------------------------------------------

El modelo es una U-Net~\cite{ronneberger2015} con codificador
ResNet-34~\cite{he2016} preentrenado en ImageNet, de la biblioteca
\emph{Segmentation Models PyTorch}~\cite{iakubovskii2019}. Produce un
mapa de probabilidades de una sola clase (filamento) a $512\times512$.

Salvo en el ajuste de la Sección~\ref{sec:ajuste-hp}, todas las
ejecuciones usan la misma configuración: 30 épocas, lotes de 8,
optimizador AdamW~\cite{loshchilov2019} con tasa de aprendizaje
$3\cdot10^{-4}$ y \emph{weight decay} $10^{-4}$, y tasa de aprendizaje
con decaimiento coseno. La red se ejecuta en precisión mixta, pero las
pérdidas se calculan siempre en \texttt{float32}: en \texttt{float16}
las sumas de clDice sobre un lote (unos $2\cdot10^6$ píxeles) desbordan.
La semilla es 42; la baseline y Topograph se repiten además con la
semilla 43.

\paragraph{Selección de la época.} Al final de cada época se evalúa el
modelo en $A$ y se guarda el mejor. El criterio es fijo para todas las
versiones: PQ global a $512\times512$, con umbral $0{,}5$ y sin filtro
de área. Se usa un post-proceso estricto a propósito: si la época se
eligiera con el post-proceso de cada versión, la comparación entre
pérdidas mezclaría dos efectos.

% ------------------------------------------------------------
\section{Pérdidas}
\label{sec:perdidas}
% ------------------------------------------------------------

Todas las pérdidas incluyen BCE, con peso $0{,}5$, y se obtienen
cambiando una sola pieza respecto a la baseline
(Tabla~\ref{tab:perdidas}). IoU se definió en la
Sección~\ref{sec:competicion}, Dice y clDice en la
Sección~\ref{sec:dice-cldice} y Topograph en la
Sección~\ref{sec:topograph}; aquí solo se fijan los pesos.

\begin{table}[ht]
\centering
\begin{tabular}{@{}ll@{}}
\toprule
Nombre & Pérdida\\
\midrule
Baseline & $0{,}5\,\mathrm{BCE}+0{,}5\,\mathrm{Dice}$\\
IoU & $0{,}5\,\mathrm{BCE}+0{,}5\,\mathrm{IoU}$\\
clDice & $0{,}5\,\mathrm{BCE}+0{,}25\,\mathrm{Dice}+0{,}25\,\mathrm{clDice}$\\
IoU+clDice & $0{,}5\,\mathrm{BCE}+0{,}25\,\mathrm{IoU}+0{,}25\,\mathrm{clDice}$\\
clDice sin Dice & $0{,}5\,\mathrm{BCE}+0{,}5\,\mathrm{clDice}$\\
\midrule
\dots{} + Topograph & la anterior $+\,\alpha\,\mathcal{L}_{\mathrm{Topo}}$, $\alpha=0{,}04$\\
\bottomrule
\end{tabular}
\caption{Pérdidas comparadas. Cada una de las cinco primeras se entrena
también con el término de Topograph.}
\label{tab:perdidas}
\end{table}

\paragraph{clDice.} Se usa el esqueleto blando de Shit et
al.~\cite{shit2021}, con 10 iteraciones de erosión.

\paragraph{Topograph.} Se usa la implementación oficial~\cite{lux2025},
con 8-conectividad, la misma que el post-proceso. Dos decisiones fijan
cómo se añade:
\begin{itemize}
  \item \emph{Peso.} Con un modelo ya entrenado, el término de
  Topograph es unas 27 veces mayor que BCE+Dice. Con $\alpha=0{,}04$
  ambos pesan lo mismo. El valor se fija así y no se ajusta en las
  comparaciones de pérdidas.
  \item \emph{Calentamiento.} Con un modelo sin entrenar, la predicción
  tiene cientos de regiones críticas por imagen y el término no aporta
  una señal útil. Por eso Topograph se activa a partir de la sexta
  época; en las cinco primeras $\alpha=0$.
\end{itemize}

% ------------------------------------------------------------
\section{Post-proceso}
\label{sec:postproceso}
% ------------------------------------------------------------

Las instancias se obtienen del mapa de probabilidades $p$ en tres pasos:
\begin{enumerate}
  \item se reescala $p$ a $2048\times2048$ (interpolación bilineal);
  \item se umbraliza, $X_t=\{p>t\}$;
  \item cada componente 8-conexa de $X_t$ es una instancia, y se
  descartan las de área menor que $a$.
\end{enumerate}
En la lectura de la Sección~\ref{sec:lectura-topologica}, son las
componentes del conjunto de nivel superior $X_t$ con un filtro que solo
mira su tamaño. Se consideran dos variantes:
\begin{itemize}
  \item \textbf{B1:} solo umbral, $t$ ajustado y $a=0$;
  \item \textbf{B2:} umbral y área mínima ajustados conjuntamente.
\end{itemize}
Los parámetros se buscan en $A$ con Optuna~\cite{akiba2019} (búsqueda
bayesiana TPE, 60 intentos, semilla fija), en los rangos
$t\in[0{,}2,\,0{,}9]$ y $a\in[0,\,300]$ píxeles a $512\times512$
(el área se escala por 16 a resolución completa). Cada intento evalúa
con la métrica oficial en $A$. El post-proceso se aplica a los
\emph{logits} ya guardados de cada ejecución, sin reentrenar.

% ------------------------------------------------------------
\section{Ajuste de hiperparámetros}
\label{sec:ajuste-hp}
% ------------------------------------------------------------

Las comparaciones de pérdidas usan los mismos hiperparámetros para
todas. Puede ocurrir que una pérdida salga peor solo porque esos
valores le sientan peor. Para comprobarlo, se ajustan los
hiperparámetros de entrenamiento de la baseline y de las dos pérdidas
más relevantes para el trabajo: Topograph, que es el método topológico
central, y clDice, que apunta a las barbas.

Se usa Optuna con TPE y, para cada modelo, se buscan tres valores: la
tasa de aprendizaje en $[10^{-4},10^{-3}]$, el \emph{weight decay} en
$[10^{-6},10^{-3}]$ (ambos en escala logarítmica) y un peso propio de
cada pérdida: el de Dice en la baseline ($[0{,}25,\,0{,}75]$, con
BCE el resto hasta 1), $\alpha$ en Topograph ($[0{,}01,\,0{,}16]$) y el
de clDice ($[0{,}10,\,0{,}40]$, restado del de Dice). El primer intento de
cada búsqueda es la configuración de la Sección~\ref{sec:entrenamiento}.
El objetivo es el PQ oficial en $A$. Los intentos que, a partir de la
época 10, van por debajo de la mediana de los anteriores en esa misma
época se interrumpen. Cada modelo dispone del mismo tiempo de cómputo.
\textbf{[completar: número de intentos por modelo y GPU.]}

% ------------------------------------------------------------
\section{Evaluación}
\label{sec:evaluacion}
% ------------------------------------------------------------

\paragraph{Métrica oficial.} Se usa el código del cuaderno oficial de
autoevaluación de la competición~\cite{kaggle2026}, sin modificarlo. La
cifra principal es el PQ global en $B$ a resolución completa
(Sección~\ref{sec:competicion}), con SQ y RQ, y las distribuciones de
la rúbrica: IoU, Dice y relaciones 1:$n$ y $n$:1.

\paragraph{Medidas propias.} Como complemento, nunca en sustitución de
las oficiales, se cuentan en $B$:
\begin{itemize}
  \item \emph{partidos y unidos:} filamentos de referencia cubiertos por
  varias predicciones, más predicciones que cubren varios filamentos;
  \item \emph{sin contacto:} filamentos sin ninguna predicción que los
  toque, más predicciones que no tocan ningún filamento.
\end{itemize}

\paragraph{Barbas.} PQ no distingue si un acierto incluye las barbas.
Para medirlo se usa la espina que MAGFiLO anota en cada
filamento~\cite{ahmadzadeh2024}. Sea $w$ el grosor medio del filamento
(área entre longitud de la espina). Las \emph{barbas} son la parte del
esqueleto de la máscara anotada que queda fuera de un tubo de radio
$1{,}5\cdot w/2$ alrededor de la espina; se descartan trozos de menos de
6 píxeles. La \emph{zona de barbas} es el entorno del filamento, hasta
$2w$, sin ese tubo ni los demás filamentos. Sobre ella se definen:
\begin{itemize}
  \item \emph{cobertura:} fracción de los píxeles de barbas reales que
  cubre la predicción;
  \item \emph{precisión:} fracción del esqueleto de la predicción en la
  zona de barbas que cae sobre una barba real;
  \item su $F_1$ y el número de ramas inventadas por filamento.
\end{itemize}
Todo se mide a $2048\times2048$, por anotación de $B$. Es una medida
propia de este trabajo.

\paragraph{Incertidumbre.} $B$ tiene 66 días, y una diferencia de PQ
puede deberse a qué días han caído en él. Para estimar esa variabilidad
se usa un \emph{bootstrap} pareado por días~\cite{efron1993}: se
remuestrean con reemplazo los días de $B$ (2.000 veces), se recalcula el
PQ global de cada ejecución con los conteos de esos días y se toma el
intervalo del 95\,\% de la diferencia con la referencia. El mismo
remuestreo se aplica a todas las ejecuciones, para que la comparación
sea pareada. Se remuestrean días, y no pares, porque los pares de un
mismo día no son independientes.

\paragraph{Variabilidad entre entrenamientos.} El entrenamiento en GPU
no es determinista: reentrenar la misma configuración con la misma
semilla cambia el PQ en $B$ entre $0{,}01$ y $0{,}02$. Las diferencias
menores que eso no se interpretan como mejoras.

% ------------------------------------------------------------
\section{Reproducibilidad}
\label{sec:reproducibilidad}
% ------------------------------------------------------------

Cada ejecución queda identificada por la fecha, la configuración y el
\emph{commit} del código, y guarda su configuración completa, la época
elegida, las predicciones y los \emph{logits} de validación. Una tabla
única registra una fila por ejecución con sus métricas en $A$ y $B$. Los
resultados del Capítulo~\ref{cap:resultados} salen de esa tabla.
